\chapter{Images}

\section{Insérer une image}

Pour insérer une image, un package intéressant est le package \jargon{graphicx} (Ajouter la commande \verb!\usepackage{graphicx}! dans le préambule).

Il fournit la commande \verb!\includegraphics! qui possède des options facultatives.

\bigskip

\begin{minipage}{0.6\linewidth}
\color{blue}\rule{9cm}{1pt}\color{black}
\vspace*{-12pt}

\begin{small}
\begin{verbatim}
\includegraphics{exemples/image}
\end{verbatim}
\end{small}

\vspace*{-12pt}
\color{blue}\rule{9cm}{1pt}\color{black}
\end{minipage}
\hfill
\begin{minipage}{0.35\linewidth}
\includegraphics{exemples/image}
\end{minipage}

\medskip

Les formats des images acceptés par compilation \texttt{PDF}\LaTeX{} sont \jargon{.jpg}, \jargon{.png} et \jargon{.pdf}.

Le format accepté par compilation \LaTeX est \jargon{.eps}.

Il est inutile de préciser l'extension du fichier utilisé. Il faut juste utiliser un fichier dont le format est le bon, suivant la compilation choisie.

\begin{info}
Si l'on veut centrer une image, il suffit d'utiliser l'environnement \jargon{center} comme pour le texte.
\end{info}

\section{Modifier une image}

La commande plusieurs \verb!\includegraphics! admet plusieurs options que l'on peut spécifier entre crochets. Ces options permettent de transformer l'aspect de l'image.

Parmi les options possibles, on trouve :

$\star$ \verb!width! : par exemple \verb!width=1cm! fixe la largeur de l'image

On peut également utiliser \verb!width=\linewidth! ou \verb!width=\textwidth!, \verb!width=0.5\linewidth!

$\star$ \verb!height! : par exemple \verb!height=2cm! fixe la hauteur de l'image

$\star$ \verb!scale! : par exemple \verb!scale=0.5! met l'image à l'échelle 0.5

$\star$ \verb!angle! : par exemple \verb!angle=45! effectue une rotation de l'image de 45\degres{} dans le sens positif, autour du centre de l'image.

\bigskip

\begin{minipage}{0.8\linewidth}
\color{blue}\rule{9cm}{1pt}\color{black}
\vspace*{-12pt}

\begin{small}
\begin{verbatim}
\includegraphics[angle=45,width=0.6\linewidth]{exemples/image}
\end{verbatim}
\end{small}

\vspace*{-12pt}
\color{blue}\rule{9cm}{1pt}\color{black}
\end{minipage}
\hfill
\begin{minipage}{0.2\linewidth}
\includegraphics[angle=45,width=0.6\linewidth]{exemples/image}
\end{minipage}

\section{Les environnements \jargon{wrapfigure} et \jargon{minipage}}

Pour utiliser l'environnement \jargon{wrapfigure}, il faut charger, dans le préambule, le package \jargon{wrapfigure}.

\VerbatimInput[label={[Environnement wrapfigure]},gobble=0]{exemples/wrapfigure.tex}

Ce \jargon{code source} donne :

\begin{wrapfigure}{r}{5cm}
\includegraphics[width=4.5cm]{exemples/image}
\end{wrapfigure}
\LaTeX{}, c'est génial ! \LaTeX{}, c'est génial ! \LaTeX{}, c'est génial ! 
\LaTeX{}, c'est génial ! \LaTeX{}, c'est génial ! \LaTeX{}, c'est génial ! 
\LaTeX{}, c'est génial ! \LaTeX{}, c'est génial ! \LaTeX{}, c'est génial ! 
\LaTeX{}, c'est génial ! \LaTeX{}, c'est génial ! \LaTeX{}, c'est génial ! 
\LaTeX{}, c'est génial ! \LaTeX{}, c'est génial ! \LaTeX{}, c'est génial ! 
\LaTeX{}, c'est génial ! \LaTeX{}, c'est génial ! \LaTeX{}, c'est génial ! 
\LaTeX{}, c'est génial ! \LaTeX{}, c'est génial ! \LaTeX{}, c'est génial ! 
\LaTeX{}, c'est génial ! \LaTeX{}, c'est génial ! \LaTeX{}, c'est génial ! 
\LaTeX{}, c'est génial ! \LaTeX{}, c'est génial ! \LaTeX{}, c'est génial ! 
\LaTeX{}, c'est génial ! \LaTeX{}, c'est génial ! \LaTeX{}, c'est génial !
\LaTeX{}, c'est génial ! \LaTeX{}, c'est génial ! \LaTeX{}, c'est génial ! 
\LaTeX{}, c'est génial ! \LaTeX{}, c'est génial ! \LaTeX{}, c'est génial !
\LaTeX{}, c'est génial ! \LaTeX{}, c'est génial ! \LaTeX{}, c'est génial ! 
\LaTeX{}, c'est génial ! \LaTeX{}, c'est génial ! \LaTeX{}, c'est génial !


\begin{info}
L'image peut être remplacée par une figure (par exemple \jargon{tikzpicture}), du texte, \dots
\end{info}

Personnellement, j'utilise plutôt l'environnement \jargon{minipage} qui permet de horizontalement découper la page en plusieurs \og sous-pages \fg{}.

\jargon{minipage} s'utilise ainsi :

\VerbatimInput[label={[Environnement minipage]},gobble=0]{exemples/minipage.tex}

Ce \jargon{code source} donne :

\bigskip

\begin{minipage}{0.7\linewidth}
\LaTeX{}, c'est génial ! \LaTeX{}, c'est génial ! \LaTeX{}, c'est génial ! 
\LaTeX{}, c'est génial ! \LaTeX{}, c'est génial ! \LaTeX{}, c'est génial ! 
\LaTeX{}, c'est génial ! \LaTeX{}, c'est génial ! \LaTeX{}, c'est génial ! 
\LaTeX{}, c'est génial ! \LaTeX{}, c'est génial ! \LaTeX{}, c'est génial ! 
\LaTeX{}, c'est génial ! \LaTeX{}, c'est génial ! \LaTeX{}, c'est génial ! 
\LaTeX{}, c'est génial ! \LaTeX{}, c'est génial ! \LaTeX{}, c'est génial ! 
\LaTeX{}, c'est génial ! \LaTeX{}, c'est génial ! \LaTeX{}, c'est génial ! 
\LaTeX{}, c'est génial ! \LaTeX{}, c'est génial ! \LaTeX{}, c'est génial ! 
\LaTeX{}, c'est génial ! \LaTeX{}, c'est génial ! \LaTeX{}, c'est génial ! 
\LaTeX{}, c'est génial ! \LaTeX{}, c'est génial ! \LaTeX{}, c'est génial ! 
\LaTeX{}, c'est génial ! \LaTeX{}, c'est génial ! \LaTeX{}, c'est génial ! 
\LaTeX{}, c'est génial ! \LaTeX{}, c'est génial ! \LaTeX{}, c'est génial ! 
\end{minipage}
\hfill
\begin{minipage}{0.25\linewidth}
\includegraphics[width=4.5cm]{exemples/image}
\end{minipage}

\medskip

\LaTeX{}, c'est génial ! \LaTeX{}, c'est génial ! \LaTeX{}, c'est génial ! 
\LaTeX{}, c'est génial ! \LaTeX{}, c'est génial ! \LaTeX{}, c'est génial ! 
\LaTeX{}, c'est génial ! \LaTeX{}, c'est génial ! \LaTeX{}, c'est génial !

\section{Une image dans du texte}

On désire parfois insérer une image dans une ligne de texte.

On  peut, par exemple, utiliser la commande \verb!\raisebox! pour cela.

\VerbatimInput[label={[Commande raisebox]},gobble=0]{exemples/raisebox1.tex}

Ce \jargon{code source} donne :

Voici du texte et une image à l'intérieur 
sans raisebox
\includegraphics[scale=0.3]{exemples/image}
\dots et avec raisebox \raisebox{-1ex}{
\includegraphics[scale=0.3]{exemples/image}}.

Ce n'est pas toujours heureux 
typographiquement parlant.

\begin{info}
La commande \texttt{\textbackslash raisebox} n'est pas réservée aux images.
\end{info}

\VerbatimInput[label={[Commande raisebox]},gobble=0]{exemples/raisebox2.tex}

Ce \jargon{code source} donne : \hspace*{10em}
Du texte \raisebox{-1ex}{en bas}
et du texte \raisebox{1ex}{en haut}.



